
\section{Architecture Integration}

\subsection{Model Support Matrix}

\begin{table}[H]
\centering
\small
\begin{tabular}{lllll}
\toprule
\textbf{Architecture} & \textbf{Nested?} & \textbf{RoPE?} & \textbf{GQA?} & \textbf{Default path} \\
\midrule
DistilBERT & No & No & No & \textbf{Dense} (encoder) \\
RoBERTa & Yes & No & No & \textbf{Dense} (encoder) \\
GPT-2 & No & No & No & \textbf{Cascade} (decoder) \\
Qwen3.5 & No & Yes & Yes & \textbf{Cascade} (decoder) \\
\bottomrule
\end{tabular}
\end{table}

\subsection{Why Encoders Use Dense by Default}

The flag \texttt{enable\_cascade\_on\_encoders} defaults to \texttt{True} **in the config 
but the \texttt{use\_dense} decision also checks**:

\begin{lstlisting}
use\_dense = is\_cross or (not self.is\_decoder and not self.enable\_cascade\_on\_encoders)
\end{lstlisting}

If \texttt{enable\_cascade\_on\_encoders=True} (default), encoders use cascade. If 
\texttt{False}, encoders use dense. \textbf{Ambiguity in the default} is a design 
concern.

\subsection{Weight Transfer}

Standard four-architecture copy:

\begin{lstlisting}
# GPT-2
c\_w = src.c\_attn.weight.data
if is\_conv1d:
    c\_w = c\_w.t()
q\_w, k\_w, v\_w = c\_w.chunk(3, dim=0)
o\_w = src.c\_proj.weight.data.t()

# RoBERTa
q\_w = src.self.query.weight.data
k\_w = src.self.key.weight.data
v\_w = src.self.value.weight.data
o\_w = src.output.dense.weight.data

# DistilBERT
q\_w = src.q\_lin.weight.data
# ...
\end{lstlisting}

\subsection{GQA Handling}

\begin{lstlisting}
if proj\_key in ("k", "v"):
    if src\_out == out\_d:
        target\_w.data.copy\_(w\_oriented)
    elif src\_out \textgreater{} out\_d and src\_out \% out\_d == 0:
        # Mean-pool
        rep = src\_out // out\_d
        target\_w.data.copy\_(w\_oriented.view(out\_d, rep, -1).mean(dim=1))
    elif src\_out \textless{} out\_d and out\_d \% src\_out == 0:
        # Repeat
        rep = out\_d // src\_out
        target\_w.data.copy\_(w\_oriented.repeat(rep, 1))
\end{lstlisting}

This handles GQA where $H_{kv} < H$ (Qwen, LLaMA).

\subsection{Injection}

\textbf{Algorithm 2:} Inject NASC
\begin{lstlisting}
1: new $\leftarrow$ NASCAttention(config, old=old\_attn)
2: if copy\_weights fails: raise (NOT return None)
3: if   has .attention.self:
4:      .attention.self $\leftarrow$ new; is\_nested $\leftarrow$ True
5: elif   has .attention:  .attention $\leftarrow$ new
6: elif   has .self:  .self $\leftarrow$ new; is\_nested $\leftarrow$ True
7: elif   has .self\_attn:  .self\_attn $\leftarrow$ new
8: elif   has .attn:  .attn $\leftarrow$ new
9: return new
\end{lstlisting}

\textbf{Important.} Unlike other modules, \texttt{inject\_nasc} \textbf{raises} on failure:

\begin{lstlisting}
except Exception as e:
    logger.error("inject\_nasc failed: \%s", e)
    raise e
\end{lstlisting}

This violates the "safe injection returning None" contract.

\subsection{Common Pitfalls}

\subsubsection{RFF Causal Path Memory}

The \texttt{KV\_cumsum} tensor in \texttt{\_global\_spectral\_rff\_attention} has shape 
$(B, H, L, 2F, d)$. For $L = 128$, this is manageable. For $L = 512$, it 
may OOM.

\subsubsection{Batch-Level Gate}

The gate is computed from \texttt{hidden\_states.mean(dim=1)}, which averages over 
all tokens. This eliminates per-token adaptivity.

\subsubsection{`softpick` Batch Fallback}

The fallback triggers when \textbf{any} row has zero denominator (batch-level), 
not per-row:

\begin{lstlisting}
zero\_mask = denom \textless{}= 1e-8  # shape (B, H, L, 1)
if zero\_mask.any():
    fallback = F.softmax(x, dim=dim)
    p = torch.where(zero\_mask, fallback, p / (denom + 1e-8))
\end{lstlisting}

The \texttt{torch.where} uses per-element fallback, but the \texttt{if} only triggers 
if any element is affected. This is slightly better than fully batch-level.

\subsubsection{The `\_logged\_once` Pattern}

NASC logs the path once per instance:

\begin{lstlisting}
if not hasattr(self, '\_logged\_once'):
    logger.debug("NASCAttention: L=\%d, ...")
    self.\_logged\_once = True
\end{lstlisting}

This is a fix from previous log spam.
\begin{lstlisting}

---
\end{lstlisting}

\begin{figure}[H]
\centering
\includegraphics[width=\linewidth]{figures/architecture.png}
\caption{High-level architectural diagram of nasc\_attention.}
\label{fig:arch_nasc-attention}
\end{figure}

